\section{The Rule of 1000}\label{supp:rule1000}
This section works through the second conclusion of Section~\ref{sec:horizons-ceiling} of the main paper. By \autoref{prop:saturation}(d), in continuously compounded rates, even with a 50-point margin intercept and 20\% market growth, a Rule of 1000 requires demand pull $\lambda h$ between \LamThouLo{} and \LamThouHi{} for $\alpha$ between 0.45 and 0.25. The upper-quartile growth surplus of the 2023--25 top tail (Section~\ref{sec:horizons-ranges}), at the upper bound on growth spend, corresponds to $\lambda h$ of \LamPiLo{} to \LamPiHi{} over the same range of $\alpha$, so a Rule of 1000 needs roughly \PullMultLo{} to \PullMultHi{} times the demand pull of the best public software firms today. We use only the upper bound: under the lower bound no selling and marketing expense buys growth, so the surplus is not a return to spend and implies no demand pull.

Such pull is necessarily transient. \autoref{prop:budget} shows that every point of score above the saturated level is paid for out of a finite budget of log share headroom, within one market and at given retention, and \autoref{cor:transient} turns the budget into a bound on the number of years a firm can sustain a given score. With a margin intercept of 60 points, 20\% market growth and a share cap of 50\%, a firm that starts at 1\% of its market can score 1{,}000 in at most \TransOnePct{} year, and one that starts at a thousandth of a percent in at most \TransSmall. A new product in a new market opens a new budget (the intuition after \autoref{prop:budget}), and each such episode is bounded in the same way.
